\chapter{Toetsen van de schermeditor}

\section{Toetsen van de schermeditor}

De volgende toetscombinaties voeren handelingen uit in de schermeditor van de MEGA65.

In sommige gevallen kan een programma de equivalente PETSCII-codes weergeven om dezelfde handelingen uit te voeren. Zo kan een programma \specialkey{CTRL} + \megakey{G}, welke een beltoon voortbrengt, als \screentext{CHR\$(7)} weergeven. Gebruik voor een \specialkey{ESC}-sequentie \screentext{CHR\$(27)} om de \specialkey{ESC}-toets weer te geven, gevolgd door de volgende toets van de reeks.

\section{Control-codes}
\label{appendix:controlcodes}
\index{Toetsenbord!CTRL}
\begin{center}
\setlength{\def\arraystretch{1.5}\tabcolsep}{6pt}
\begin{longtable}{c|L{5.5cm}}
	\textbf{Toetsencombinatie} & \textbf{Functie}\\
   \hhline{==}
	\endhead

  \multicolumn{2}{l}{\textbf{Kleuren}} \\
  \hhline{==}
\specialkey{CTRL} + \megakey{1} tot \megakey{8} &Kiezen uit het eerste kleurbereik. Zie bijlage \vref{appendix:colourtable} voor de lijst van kleuren in het systeempalet.\\
\hline
\megasymbolkey + \megakey{1} tot \megakey{8} &Kiezen uit het tweede kleurbereik.  \\
\hline
\specialkey{CTRL} + \megakey{E} &Herstelt de kleur van de cursor naar de standaardkleur (wit).\\
\hline
\specialkey{CTRL} + \megakey{D} &Schakelt de VIC-IV over naar kleurbereik 0-15 (standaardkleuren). Deze kleuren zijn bereikbaar met \specialkey{CTRL} en de toetsen \megakey{1} tot \megakey{8} (voor de eerste 8 kleuren) of met \megasymbolkey en de toetsen \megakey{1} tot \megakey{8} (voor de overige 8 kleuren).\\
\hline
\specialkey{CTRL} + \megakey{A} &Schakelt de VIC-IV over naar kleurbereik 16-31 (alternatieve/regenboog-kleuren). Deze kleuren zijn bereikbaar met \specialkey{CTRL} en de toetsen \megakey{1} tot \megakey{8} (voor de eerste 8 kleuren) of met \megasymbolkey en de toetsen \megakey{1} tot \megakey{8} (voor de overige 8 kleuren).\\
  \hhline{==}
  \multicolumn{2}{l}{\textbf{Tabposities}} \\
  \hhline{==}
\specialkey{CTRL} + \megakey{Z} &Verplaatst de cursor via een tab naar links. Indien er geen tabposities meer resteren, blijft de cursor aan het begin van de regel staan.\\
\hline
\specialkey{CTRL} + \megakey{I} &Verplaatst de cursor via een tab naar rechts. Indien er geen tabposities meer resteren, blijft de cursor aan het einde van de regel staan.\\
\hline
\specialkey{CTRL} + \megakey{X} &Stelt de huidige schermkolom als tabpositie in of wist deze.
 Gebruik \specialkey{CTRL} +\megakey{Z} en \megakey{I}  om heen en weer te springen naar alle posities die met \megakey{X} zijn ingesteld.\\
  \hhline{==}
  \multicolumn{2}{l}{\textbf{Verplaatsen}} \\
  \hhline{==}
\specialkey{CTRL} + \megakey{Q} &Verplaatst de cursor telkens één regel omlaag. Gelijk aan \megakey{$\downarrow$}.\\
\hline
\specialkey{CTRL} + \megakey{J} &Verplaatst de cursor één positie omlaag. Indien u zich op een lange regel BASIC-code bevindt die over twee regels doorloopt, verplaatst de cursor zich twee rijen omlaag om op de volgende regel te komen.\\
\hline
\specialkey{CTRL} + \megakey{]} &Gelijk aan \megakey{$\rightarrow$}.\\
\hline
\specialkey{CTRL} + \megakey{T} &Wist het teken direct links van de cursor en schuift alle tekens die rechts daarvan staan één positie naar links. Dit is gelijk aan \specialkey{INST\\DEL}.\\
\hline
\specialkey{CTRL} + \megakey{M} &Voert een regelterugloop uit, gelijk aan \specialkey{RETURN}.\\

  \hhline{==}
  \multicolumn{2}{l}{\textbf{Verplaatsen per woord}} \\
  \hhline{==}

\specialkey{CTRL} + \megakey{U} &Verplaatst de cursor terug naar het begin van het vorige woord. Indien er op de huidige regel geen vorig woord staat, verplaatst de cursor zich naar de eerste kolom van de huidige regel en vervolgens naar de vorige regel, totdat een regel met een woord wordt aangetroffen.\\
\hline
\specialkey{CTRL} + \megakey{W} &Verplaatst de cursor vooruit naar het begin van het volgende woord. Indien er op de huidige regel geen volgend woord staat, verplaatst de cursor zich naar de eerste kolom van de volgende regel, totdat een regel met een woord wordt aangetroffen.\\

  \hhline{==}
  \multicolumn{2}{l}{\textbf{Scrollen}} \\
  \hhline{==}

\specialkey{CTRL} + \megakey{P} &Schuift de BASIC-listing één regel omlaag. Gelijk aan \megakey{F9}.\\
\hline
\specialkey{CTRL} + \megakey{V} &Schuift de BASIC-listing één regel omhoog. Gelijk aan \megakey{F11}.\\
\hline
\specialkey{CTRL} + \megakey{S} &Gelijk aan \specialkey{NO\\SCROLL}.\\

  \hhline{==}
  \multicolumn{2}{l}{\textbf{Opmaak}} \\
  \hhline{==}

\specialkey{CTRL} + \megakey{B} &Schakelt de onderstrepingsmodus voor tekst in. U kunt de onderstrepingsmodus uitschakelen door \specialkey{ESC} en vervolgens \megakey{O} in te drukken.\\
\hline
\specialkey{CTRL} + \megakey{O} &Schakelt de flitsmodus voor tekst in. U kunt de flitsmodus uitschakelen door \specialkey{ESC} en vervolgens \megakey{O} in te drukken.\\
\hline

  \hhline{==}
  \multicolumn{2}{l}{\textbf{Hoofdlettergebruik}} \\
  \hhline{==}

\specialkey{CTRL} + \megakey{N} &Wijzigt het hoofdlettergebruik van de tekst van uitsluitend hoofdletters naar kleine letters.\\
\hline
\specialkey{CTRL} + \megakey{K} &Vergrendelt de omschakeling tussen hoofdletters en kleine letters, welke gewoonlijk met \megasymbolkey + \specialkey{SHIFT} wordt uitgevoerd.\\
\hline
\specialkey{CTRL} + \megakey{L} &Schakelt de omschakeling tussen hoofdletters en kleine letters in, welke met \megasymbolkey + \specialkey{SHIFT} wordt uitgevoerd.\\

  \hhline{==}
  \multicolumn{2}{l}{\textbf{Diversen}} \\
  \hhline{==}

\specialkey{CTRL} + \megakey{G} &Brengt een beltoon voort.\\
\hline
\specialkey{CTRL} + \megakey{[} &Gelijk aan het indrukken van \specialkey{ESC}.\\
\hline
\specialkey{CTRL} + \megakey{*} &Opent de Matrix Mode Debugger.\\
\hline

\end{longtable}
\end{center}


\section{Shiftcodes}
\label{appendix:shiftedcodes}
\index{Toetsenbord!Shifttoetsen}
\begin{center}
\setlength{\def\arraystretch{1.5}\tabcolsep}{6pt}
\begin{longtable}{c|L{5.5cm}}
	\textbf{Toetsencombinatie} & \textbf{Functie}\\
  \hhline{==}
	\endhead

\specialkey{SHIFT} + \specialkey{INST\\DEL} &Voegt op de huidige cursorpositie een teken in en verplaatst alle tekens één positie naar rechts.\\
\hline
\specialkey{SHIFT} + \specialkey{HOME} &Wist het volledige scherm en plaatst de cursor op de startpositie.\\
\hline

\end{longtable}
\end{center}



\section{Escape-sequenties}
\label{escape-sequences}
\index{Toetsenbord!Escape-reeksen}
Om een escape-sequentie uit te voeren, drukt u \specialkey{ESC} kort in en laat u deze weer los. Druk vervolgens één van de volgende toetsen in om de sequentie uit te voeren.

\begin{center}
\setlength{\def\arraystretch{1.5}\tabcolsep}{6pt}
\begin{longtable}{c|L{5.5cm}}
	\textbf{Toets} & \textbf{Sequentie}\\
  \hhline{==}
	\endhead

  \multicolumn{2}{l}{\textbf{Gedrag van de editor}} \\
  \hhline{==}

\specialkey{ESC} \megakey{X} &
Wist het scherm en schakelt tussen de tekstmodi 40 $\times$ 25 en 80 $\times$ 25.\\
\hline

\specialkey{ESC} \megakey{4} &
Wist het scherm en schakelt over naar de tekstmodus 40 $\times$ 25.\\
\hline

\specialkey{ESC} \megakey{8} &
Wist het scherm en schakelt over naar de tekstmodus 80 $\times$ 25.\\
\hline

\specialkey{ESC} \megakey{5} &
Schakelt over naar de tekstmodus 80 $\times$ 50.\\
\hline
\multicolumn{2}{L{7cm}}{  Let op dat sommige programma's verwachten dat zij in de tekstmodus 80 $\times$ 25 worden gestart en zich mogelijk niet correct gedragen wanneer zij in de tekstmodus 80 $\times$ 50 worden gestart.} \\
\hline

\specialkey{ESC} \megakey{@} &
%TODO the megakey with @ shows the label far smaller than any other key.
% Perhaps it needs more styling?
Wist een gedeelte van het scherm, vanaf de huidige cursorpositie tot het einde van het scherm.\\
\hline

\specialkey{ESC} \megakey{O} &
Schakelt de zinnebeeldmodus, de modus voor omgekeerde tekst, de onderstrepingsmodus en de flitsmodus uit.\\

  \hhline{==}
  \multicolumn{2}{l}{\textbf{Scrollen}} \\
  \hhline{==}
% TODO for some reason this doesn't leave enough space (the keys overlap with the hhline)
\specialkey{ESC} \megakey{V} &
Schuift het hele scherm één regel omhoog.\\
\hline
\specialkey{ESC} \megakey{W} &
Schuift het hele scherm één regel omlaag.\\
\hline
\specialkey{ESC} \megakey{L} &
Schakelt het scrollen in wanneer \megakey{$\downarrow$} onderaan het scherm wordt ingedrukt.\\
\hline
\specialkey{ESC} \megakey{M} &
Schakelt het scrollen uit. Wanneer \megakey{$\downarrow$} onderaan het scherm wordt ingedrukt, verplaatst de cursor zich naar de bovenkant van het scherm. Wanneer \megakey{$\uparrow$} echter bovenaan het scherm wordt ingedrukt, blijft de cursor op de eerste regel staan.\\
\hline
\specialkey{ESC} \megakey{N} &
Schakelt het ``regelduwen'' in: typen of afdrukken in de meest rechtse kolom schuift de daaropvolgende regels één positie omlaag.\\
\hline
\specialkey{ESC} \megakey{R} &
Schakelt het ``regelduwen'' uit: typen of afdrukken in de meest rechtse kolom verplaatst de cursor naar het begin van de volgende regel, maar schuift geen regels omlaag. Schakel zowel het ``regelduwen'' (\specialkey{ESC} \megakey{R}) als het scrollen (\specialkey{ESC} \megakey{M}) uit om met {\bf PRINT} in de meest rechtse kolom af te drukken zonder de rest van de weergave te verstoren.\\
\hhline{==}
  \multicolumn{2}{l}{\textbf{Invoegen en verwijderen}} \\
  \hhline{==}
\specialkey{ESC} \megakey{I} &
Voegt een lege regel in op de huidige cursorpositie en schuift alle daaropvolgende regels één positie omlaag.\\
\hline
\specialkey{ESC} \megakey{D} &
Verwijdert de huidige regel en schuift de regels onder de cursor één positie omhoog.\\
\hline
\specialkey{ESC} \megakey{P} &
Wist alle tekens van de cursor tot aan het begin van de huidige regel.\\
\hline
\specialkey{ESC} \megakey{Q} &
Wist alle tekens van de cursor tot aan het einde van de huidige regel.\\
  \hhline{==}
  \multicolumn{2}{l}{\textbf{Verplaatsen}} \\
  \hhline{==}
\specialkey{ESC} \megakey{J} &
Verplaatst de cursor naar het begin van de huidige regel.\\
\hline
\specialkey{ESC} \megakey{K} &
Verplaatst de cursor naar het laatste niet-witruimteteken op de huidige regel.\\
\hline
\specialkey{ESC} \megakeywhite{$\uparrow$} &
\index{Toetsenbord!Cursortoetsen}
Slaat de huidige cursorpositie op. Gebruik \specialkey{ESC} \megakeywhite{$\leftarrow$} (naast \megakey{1}) om de cursor terug te verplaatsen naar de opgeslagen positie. Let op: de \megakeywhite{$\uparrow$} die hier gebruikt wordt, staat naast \widekey{RESTORE}.\\
\hline
\specialkey{ESC} \megakeywhite{$\leftarrow$} &
Herstelt de cursorpositie naar de positie die eerder is opgeslagen door het indrukken van de toets \specialkey{ESC}
    \megakeywhite{$\uparrow$}
(naast \widekey{RESTORE}). Let op: de \megakeywhite{$\leftarrow$} die hier gebruikt wordt, staat naast \megakey{1}.\\
\hline
\specialkey{ESC} \specialkey{HOME} &
Herstelt de cursorpositie naar de positie die eerder is opgeslagen door het indrukken van \specialkey{HOME}.\\

  \hhline{==}
  \multicolumn{2}{l}{\textbf{Vensters}} \\
  \hhline{==}
\specialkey{ESC} \megakey{T} &
Stelt de linkerbovenhoek van het venstergebied in. Alle ingevoerde tekens en schermactiviteit blijven tot dit gebied beperkt. Zie ook \specialkey{ESC} \megakey{B}. U kunt de venstermodus uitschakelen door \specialkey{CLR\\HOME} tweemaal in te drukken. \\
\hline
\specialkey{ESC} \megakey{B} &
Stelt de rechterbenedenhoek van het venstergebied in. Alle ingevoerde tekens en schermactiviteit blijven tot dit gebied beperkt. Zie ook \specialkey{ESC} \megakey{T}. U kunt de venstermodus uitschakelen door \specialkey{CLR\\HOME} tweemaal in te drukken.\\
  \hhline{==}
  \multicolumn{2}{l}{\textbf{Gedrag van de cursor}} \\
  \hhline{==}
\specialkey{ESC} \megakey{A} &
Schakelt de automatische invoegmodus in. Elke ingedrukte toets wordt op de huidige cursorpositie ingevoegd, waarna alle tekens op de huidige regel rechts van de cursor één positie naar rechts opschuiven.\\
\hline
\specialkey{ESC} \megakey{C} &
Schakelt de automatische invoegmodus uit en keert terug naar de overschrijfmodus.\\
\hline
\specialkey{ESC} \megakey{E} &
Schakelt de cursor over naar de niet-flitsende modus.\\
\hline
\specialkey{ESC} \megakey{F} &
Schakelt de cursor over naar de gewone flitsmodus.\\
  \hhline{==}
  \multicolumn{2}{l}{\textbf{Gedrag van de bel}} \\
  \hhline{==}
\specialkey{ESC} \megakey{G} &
Schakelt de bel in. U kunt de bel laten klinken met \specialkey{CTRL} en \megakey{G}.\\
\hline
\specialkey{ESC} \megakey{H} &
Schakelt de bel uit, zodat het indrukken van \specialkey{CTRL} en \megakey{G} geen effect heeft.\\
  \hhline{==}
  \multicolumn{2}{l}{\textbf{Kleuren}} \\
  \hhline{==}
\specialkey{ESC} \megakey{U} &
\label{appendix:escape-colours}
Schakelt de VIC-IV over naar kleurbereik 0-15 (standaardkleuren). Deze kleuren zijn bereikbaar met \specialkey{CTRL} en de toetsen \megakey{1} tot \megakey{8} (voor de eerste 8 kleuren) of met \megasymbolkey en de toetsen \megakey{1} tot \megakey{8} (voor de overige kleuren).\\
\hline
\specialkey{ESC} \megakey{S} &
Schakelt de VIC-IV over naar kleurbereik 16-31 (alternatieve/rainbow-kleuren). Deze kleuren zijn bereikbaar met \specialkey{CTRL} en de toetsen \megakey{1} tot \megakey{8} (voor de eerste 8 kleuren) of met \megasymbolkey en de toetsen \megakey{1} tot \megakey{8} (voor de overige kleuren).\\
\hline
  \hhline{==}
  \multicolumn{2}{l}{\textbf{Tabposities}} \\
  \hhline{==}
\specialkey{ESC} \megakey{Y} &
Stelt de standaardtabposities in (om de 8 spaties) voor het volledige scherm.\\
\hline
\specialkey{ESC} \megakey{Z} &
Wist alle tabposities. Met \specialkey{CTRL} en \megakey{I} verplaatst een tab de cursor naar het einde van de regel.\\
\hline
\end{longtable}
\end{center}
